\section{Free Abelian Groups}

\exercise{38.1}{A basis of $\Z^3$ with no zero coordinates}
Take $u=(1,1,1)$, $v=(1,2,2)$, $w=(1,2,3)$.
For any $(e,f,g)\in\Z^3$, solving $xu+yv+zw=(e,f,g)$ by subtracting coordinates gives
\[
z=g-f,\qquad y=2f-e-g,\qquad x=2e-f.
\]
These are integers and substitution recovers all three coordinates. The formulas also give a unique coefficient triple, so the vectors generate $\Z^3$ and are independent over $\Z$. Thus they form a basis, with every coordinate nonzero.

\exercise{38.2}{The vectors $(2,1)$ and $(3,1)$}
Yes. For an arbitrary $(e,f)$, the equations $2x+3y=e$, $x+y=f$ give
\[
y=e-2f,\qquad x=3f-e.
\]
Both coefficients are integers, so every integer pair is generated. They are uniquely determined, so the only representation of zero is $x=y=0$. These are precisely the generating and independence requirements for a basis of a free abelian group.

\exercise{38.3}{The vectors $(2,1)$ and $(4,1)$}
No. Any integer combination $x(2,1)+y(4,1)=(2x+4y,x+y)$ has even first coordinate. Therefore $(1,0)$ is not generated.
The two vectors are linearly independent over $\R$ (their determinant is $-2$), but that does not make them a $\Z$-basis: integer coefficients, rather than arbitrary real coefficients, are required.

\exercise{38.4}{The determinant criterion for an integer basis}
Put $D=ad-bc$, the determinant of the matrix with columns $(a,b)$ and $(c,d)$.
If $D=\pm1$, solving $x(a,b)+y(c,d)=(e,f)$ gives
\[
x=\frac{de-cf}{D},\qquad y=\frac{af-be}{D}.
\]
These coefficients are integers and unique for all $e,f\in\Z$, proving the basis property.
Conversely, if the two vectors are a basis, there is an integer matrix $B$ whose columns express $(1,0),(0,1)$ in that basis. For $A=\begin{pmatrix}a&c\\b&d\end{pmatrix}$ this gives $AB=I$.
Taking determinants yields $D\det B=1$. Both factors are integers, so $D=\pm1$.
Hence the necessary and sufficient condition is $ad-bc=\pm1$.

\exercise{38.5}{Definition of rank}
Replace ``a generating set'' by ``a basis.'' The rank is the cardinality of any basis of the free abelian group. Generating sets can contain redundant elements: $\{1\}$ and $\{1,2\}$ both generate $\Z$ but have different sizes. Bases have invariant cardinality.

\exercise{38.6}{Definition of an abelian basis}
The proposed definition is correct. The zero relation with distinct basis elements forces every integer coefficient to vanish, and the set generates the group. The reverse implication in the independence condition is automatic, since a sum with all coefficients zero is zero. A group having such a basis is free abelian.

\exercise{38.7}{A proper subgroup of the same rank}
For any positive integer $r$, take $G=\Z^r$ and $K=2\Z\times\Z^{r-1}$. It is proper since $(1,0,\ldots,0)\notin K$.
The $r$ vectors $2e_1,e_2,\ldots,e_r$ generate $K$ and are independent over $\Z$, so $K$ has rank $r$, just as $G$ does. In particular already $2\Z<\Z$ gives a rank-1 example.

\exercise{38.8}{True or false: free abelian groups}
\begin{enumerate}[label=\textbf{(\alph*)}]
\item True. If a multiple of an element is zero, uniqueness of basis coefficients makes every coefficient zero.
\item True by the structure theorem: the torsion cyclic factors vanish, leaving $\Z^r$.
\item True: $\Z^r$ has the standard $r$-element basis.
\item True. If $G$ has Betti number $r$ and is generated by $r$ elements, write $G\cong\Z^r\oplus T$. A prime $p\mid|T|$ would give $G/pG\cong\mathbb F_p^r\oplus T/pT$ of dimension greater than $r$, impossible with $r$ generators. Thus $T=0$.
\item True. The larger set still contains every generator in $X$, so still generates all of $G$.
\item False. $\{1\}$ is a basis of $\Z$, but $\{1,2\}$ is dependent: $2\cdot1-1\cdot2=0$.
\item False. $\Z$ has only two bases, $\{1\}$ and $\{-1\}$.
\item True. If $x_1,x_2,\ldots$ is a basis, replacing $x_1$ by $x_1+kx_2$ gives a different basis for each $k\in\Z$, since the old $x_1$ is recovered by subtracting $kx_2$.
\item True: a subgroup of $\Z^r$ is free abelian, as proved in this section.
\item False. For $G=\Z$, $K=2\Z$, the quotient is $\Z_2$, which has nonzero torsion and is not free abelian.
\end{enumerate}

\exercise{38.9}{The coordinate isomorphism}
For a basis $x_1,\ldots,x_r$, define $\phi(\sum n_ix_i)=(n_1,\ldots,n_r)$.
Uniqueness of a basis expansion makes this well defined, including the zero vector. If $a=\sum n_ix_i$ and $b=\sum m_ix_i$, then $a+b=\sum(n_i+m_i)x_i$, so $\phi(a+b)=\phi(a)+\phi(b)$.
If $\phi(a)=0$, all coefficients are zero, hence $a=0$; thus the map is injective. Given any integer tuple $(n_1,\ldots,n_r)$, its preimage is $\sum n_ix_i$, so the map is onto. This completes the claimed isomorphism $G\cong\Z^r$.

\exercise{38.10}{A free abelian group has no nonzero torsion}
Let $a=\sum_{i=1}^r n_ix_i$ in a basis, with finitely many coefficients involved, and suppose $ma=0$ for an integer $m\ne0$. Then $\sum_i(mn_i)x_i=0$.
Independence forces $mn_i=0$ in $\Z$ for each $i$. Since $m\ne0$, every $n_i=0$, so $a=0$. Thus a nonzero element cannot have finite order; the argument works even for an infinite basis because each expansion is finite.

\exercise{38.11}{Products of free abelian groups}
If $X$ is a basis of $G$ and $Y$ a basis of $G'$, use
\[
\{(x,0):x\in X\}\ \cup\ \{(0,y):y\in Y\}.
\]
It generates $G\times G'$ by expanding each coordinate separately. A zero integer combination has zero first coordinate, forcing all $X$-coefficients to be zero, and zero second coordinate, forcing all $Y$-coefficients to be zero. Thus the displayed set is a basis, and the product is free abelian. This is a finite product of groups; even when the bases are infinite, each individual expansion remains finite.

\exercise{38.12}{Finitely generated torsion-free groups}
A free abelian group of finite rank has a finite basis and hence is finitely generated; Exercise 38.10 proves it is torsion free.
Conversely the structure theorem writes a finitely generated abelian group as $\Z^r\oplus T$, where $T$ is finite. If there is no nonzero element of finite order, $T$ must be trivial. Thus the group is $\Z^r$, which has its standard basis.
The zero group is included by the usual empty-basis/rank-zero convention; if ``free abelian'' is restricted to nonzero groups as in the section's initial definition, state the equivalence for nonzero groups.

\exercise{38.13}{Why $\Q$ is not free abelian}
Two nonzero rational numbers $x=a/b$ and $y=c/d$ satisfy
$(bc)x-(ad)y=0$, with both coefficients nonzero when $a,c\ne0$. Thus no independent set over $\Z$ can contain two distinct nonzero rational numbers; zero cannot be a basis element either.
A basis for the nonzero group $\Q$ would therefore contain just one nonzero rational $q$. But its generated subgroup consists of integer multiples of $q$, and $q/2$ is not one of them. Hence one element cannot generate $\Q$, so $\Q$ is not free abelian.

\exercise{38.14}{The $p$-primary torsion subgroup}
The set $T_p$ contains zero. If $p^a x=0$ and $p^b y=0$, let $c=\max(a,b)$; then $p^c(x-y)=0$. Hence $x-y$ has order a power of $p$ (or is zero), and lies in $T_p$.
The one-step subgroup criterion proves $T_p\le T$. The commutativity of $T$ is used when distributing multiplication by $p^c$ over a difference.

\exercise{38.15}{Locating the $p$-primary factors}
In a prime-power product decomposition of $T$, an element is killed by a power of $p$ exactly when all its coordinates in $q$-power factors with $q\ne p$ vanish. Indeed, multiplication by $p$ is invertible on a cyclic $q$-group, because $\gcd(p,q)=1$. Its repeated powers therefore cannot kill a nonzero coordinate.
Conversely the remaining $p$-power coordinates are all killed by some common power of $p$. Restricting the product isomorphism to $T_p$ thus identifies it exactly with the product of the cyclic $p$-power factors.

\exercise{38.16}{Elements killed by an integer}
The multiplication map $\mu_n:G\to G$, $x\mapsto nx$, is a homomorphism because $G$ is abelian. Therefore $G[n]=\ker\mu_n$ is a subgroup.
Directly, $n0=0$, and if $nx=ny=0$, then $n(x-y)=nx-ny=0$, proving the same fact without invoking kernels. For a general nonabelian group the power map need not be a homomorphism; the abelian assumption is essential.

\exercise{38.17}{The $p$-torsion of a cyclic $p$-group}
In $\Z_{p^r}$ the equation $px=0$ means $p^r\mid px$, equivalently $p^{r-1}\mid x$. Thus
\[
\Z_{p^r}[p]=\{0,p^{r-1},2p^{r-1},\ldots,(p-1)p^{r-1}\}.
\]
It is cyclic generated by $p^{r-1}$ and has order $p$, so the map $\Z_p\to\Z_{p^r}[p]$ sending $k$ to $kp^{r-1}$ is a well-defined bijective homomorphism.

\exercise{38.18}{Counting $p$-torsion in a product}
An element of a direct product is killed by $p$ exactly when every coordinate is killed by $p$. Hence
\[
\left(\prod_{i=1}^m\Z_{p^{r_i}}\right)[p]
=\prod_{i=1}^m\bigl(\Z_{p^{r_i}}[p]\bigr)
\cong\Z_p^m
\]
by Exercise 38.17. In particular its size is $p^m$; this size is intrinsic to the group, independent of a chosen decomposition.

\exercise{38.19}{Uniqueness of the $p$-power exponents}
\textbf{(a)} Apply the intrinsic construction $T_p[p]$ to the two decompositions. Their orders are $p^m$ and $p^n$, so $m=n$.
\textbf{(b)} Arrange exponents increasingly. Suppose that $r_i=s_i$ for $i<j$ and, seeking a contradiction, $r_j<s_j$. The group $p^{r_j}T_p$ has from the first decomposition one nonzero cyclic factor for each $r_i>r_j$, at most $m-j$ factors. From the second decomposition every index $i\ge j$ has $s_i\ge s_j>r_j$, giving at least $m-j+1$ nonzero factors; the indices before $j$ have the same behavior in both decompositions.
Part (a), applied to this intrinsic subgroup, forbids different numbers of nonzero $p$-power cyclic factors. Thus $r_j<s_j$ is impossible. Interchange the decompositions to rule out $s_j<r_j$. Start at $j=1$ and induct to obtain $r_i=s_i$ for all $i$.

\exercise{38.20}{From invariant factors to prime powers}
Factor each torsion coefficient $m_i$ into pairwise coprime prime powers. The Chinese remainder isomorphism gives
\[
\Z_{m_i}\cong\prod_{p\mid m_i}\Z_{p^{v_p(m_i)}}.
\]
Taking the product over $i$ and regrouping by prime produces the prime-power decomposition of $T$. Its uniqueness follows from Exercises 38.14--19. Factors of order 1 are omitted, since they contribute only a trivial coordinate.

\exercise{38.21}{Recovering the largest torsion coefficient}
Because $m_i\mid m_{i+1}$, the valuations $v_p(m_i)$ are nondecreasing for each prime $p$. The largest exponent for $p$ in the entire prime-power decomposition is consequently $v_p(m_r)$.
The same intrinsic largest exponent is $v_p(n_s)$ in the other decomposition. Therefore, writing it as $h_p$,
\[
m_r=\prod_{p\mid|T|}p^{h_p}=n_s.
\]
All primes dividing $T$ occur in this product, so equality of the prime valuations proves equality of the integers.

\exercise{38.22}{Recovering all torsion coefficients}
For each prime, list its exponents in the unique prime-power decomposition in increasing order. Pad each prime's list on the \emph{left} with zeros to a common length. The last torsion coefficient is the product of the prime powers in the last column; the preceding coefficient is the product in the next-to-last column, and so on.
This reconstructs $m_{r-1}$ and $n_{s-1}$ identically: for each prime take its second-largest exponent, or zero if it appears only once. Continue using third-largest exponents and so forth. Divisibility forces precisely this alignment of the lists, because valuations are nondecreasing.
The number of columns is the largest number of nonzero exponents in any prime's list: otherwise the first coefficient would be 1, contrary to $m_1,n_1>1$. Thus the common length is both $r$ and $s$, and every coefficient agrees. This proves $r=s$ and $m_i=n_i$ for all $i$.
